\input{../doc-class-cours.tex}

\begin{document}

\begin{center}
  \textit{Il faut apprendre, non pas pour l'amour de la connaissance, mais pour se défendre contre le mépris dans lequel le monde tient les ignorants.} - \textbf{Charlie Chaplin}
\end{center}

\begin{multicols}{2}
  
\subsection*{Ex1 - Nommer les angles} 

\begin{figure}[H]
  \centering
  \includegraphics[width=0.8\linewidth]{5x3-triangles/ex1a.pdf}
\end{figure} \columnbreak

\subsection*{Ex3 - Calculer des angles dans un triangle} 

\begin{figure}[H]
  \centering
  \includegraphics[width=\linewidth]{5x3-triangles/ex3a.pdf}
\end{figure} 

\end{multicols}

\horrule{1px}

\begin{multicols}{2}

\subsection*{Ex2 - Calculer des angles} 

\begin{figure}[H]
  \centering
  \includegraphics[width=0.8\linewidth]{5x3-triangles/ex2a.pdf}
\end{figure} \columnbreak

\subsection*{Ex4 - Calculer des aires}

\begin{figure}[H]
  \centering
  \includegraphics[width=\linewidth]{5x3-triangles/ex4a.pdf}
\end{figure}

\end{multicols}
\horrule{1px}
\begin{multicols}{2}

\begin{figure}[H]
  \centering
  \includegraphics[width=0.4\linewidth]{5x3-triangles/demo.pdf}
\end{figure} \columnbreak

\subsection*{PB1 - Démo} 

Démontrer que la somme des angles dans un triangle fait 180°. 

\end{multicols}
\horrule{1px}
\begin{multicols}{2}

\subsection*{PB2 - Calculer un angle} 

\begin{figure}[H]
  \centering
  \includegraphics[width=0.8\linewidth]{5x3-triangles/pb2a.pdf}
\end{figure} \columnbreak

\subsection*{PB3 - Calculer une aire} 

\begin{figure}[H]
  \centering
  \includegraphics[width=\linewidth]{5x3-triangles/pb3a.pdf}
\end{figure}

\end{multicols}

\newpage

\begin{center}
  \textit{Il faut apprendre, non pas pour l'amour de la connaissance, mais pour se défendre contre le mépris dans lequel le monde tient les ignorants.} - \textbf{Charlie Chaplin}
\end{center}

\begin{multicols}{2}
  
\subsection*{Ex1 - Nommer les angles} 

\begin{figure}[H]
  \centering
  \includegraphics[width=0.8\linewidth]{5x3-triangles/ex1b.pdf}
\end{figure} \columnbreak

\subsection*{Ex3 - Calculer des angles dans un triangle} 

\begin{figure}[H]
  \centering
  \includegraphics[width=\linewidth]{5x3-triangles/ex3b.pdf}
\end{figure} 

\end{multicols}

\horrule{1px}

\begin{multicols}{2}

\subsection*{Ex2 - Calculer des angles} 

\begin{figure}[H]
  \centering
  \includegraphics[width=0.8\linewidth]{5x3-triangles/ex2b.pdf}
\end{figure} \columnbreak

\subsection*{Ex4 - Calculer des aires}

\begin{figure}[H]
  \centering
  \includegraphics[width=\linewidth]{5x3-triangles/ex4b.pdf}
\end{figure}

\end{multicols}
\horrule{1px}
\begin{multicols}{2}

\begin{figure}[H]
  \centering
  \includegraphics[width=0.4\linewidth]{5x3-triangles/demo.pdf}
\end{figure} \columnbreak

\subsection*{PB1 - Démo} 

Démontrer que la somme des angles dans un triangle fait 180°. 

\end{multicols}
\horrule{1px}
\begin{multicols}{2}

\subsection*{PB2 - Calculer un angle} 

\begin{figure}[H]
  \centering
  \includegraphics[width=0.8\linewidth]{5x3-triangles/pb2b.pdf}
\end{figure} \columnbreak

\subsection*{PB3 - Calculer une aire} 

\begin{figure}[H]
  \centering
  \includegraphics[width=\linewidth]{5x3-triangles/pb3b.pdf}
\end{figure}

\end{multicols}

\newpage

\begin{multicols}{3}

\begin{figure}[H]
  \centering
  \includegraphics[width=0.9\linewidth]{5x3-triangles/demo.pdf}
\end{figure}  \horrule{1px}

\begin{figure}[H]
  \centering
  \includegraphics[width=0.9\linewidth]{5x3-triangles/demo.pdf}
\end{figure}  \horrule{1px}

\begin{figure}[H]
  \centering
  \includegraphics[width=0.9\linewidth]{5x3-triangles/demo.pdf}
\end{figure}  \horrule{1px}

\begin{figure}[H]
  \centering
  \includegraphics[width=0.9\linewidth]{5x3-triangles/demo.pdf}
\end{figure}  \horrule{1px}


\begin{figure}[H]
  \centering
  \includegraphics[width=0.9\linewidth]{5x3-triangles/demo.pdf}
\end{figure}  \horrule{1px}


\begin{figure}[H]
  \centering
  \includegraphics[width=0.9\linewidth]{5x3-triangles/demo.pdf}
\end{figure}  \horrule{1px}

\begin{figure}[H]
  \centering
  \includegraphics[width=0.9\linewidth]{5x3-triangles/demo.pdf}
\end{figure}  \horrule{1px}

\begin{figure}[H]
  \centering
  \includegraphics[width=0.9\linewidth]{5x3-triangles/demo.pdf}
\end{figure}  \horrule{1px}

\begin{figure}[H]
  \centering
  \includegraphics[width=0.9\linewidth]{5x3-triangles/demo.pdf}
\end{figure}  \horrule{1px}

\begin{figure}[H]
  \centering
  \includegraphics[width=0.9\linewidth]{5x3-triangles/demo.pdf}
\end{figure}  \horrule{1px}

\begin{figure}[H]
  \centering
  \includegraphics[width=0.9\linewidth]{5x3-triangles/demo.pdf}
\end{figure}  \horrule{1px}

\begin{figure}[H]
  \centering
  \includegraphics[width=0.9\linewidth]{5x3-triangles/demo.pdf}
\end{figure}  \horrule{1px}

\end{multicols}

\end{document}